\section{The \texttt{StatisticalChopper} McStas Component}
Statistical (correlation) Chopper

\subsection*{Identification}
\begin{itemize}
  \item \textbf{Author:} C. Monzat/E. Farhi/S. Rozenkranz
  \item \textbf{Origin:} ILL
  \item \textbf{Date:} Nov 10th, 2009
\end{itemize}

\subsection*{Description}
This component is a statistical chopper based on a pseudo random sequence that the user can specify. Inspired from DiskChopper, it should be used with SatisticalChopper\_Monitor component.

Example: StatisticalChopper(nu=1487*60/255) use default sequence

\subsection*{Input parameters}
Parameters in \textbf{boldface} are required; the others are optional.

\begin{longtable}{p{0.22\textwidth}p{0.12\textwidth}p{0.46\textwidth}p{0.14\textwidth}}
\toprule
\textbf{Name} & \textbf{Unit} & \textbf{Description} & \textbf{Default} \\
\midrule
\endhead
radius & m & Radius of the disc & 0.5 \\
yheight & m & Slit height (if = 0, equal to radius). Auto centering of beam at half height. & 0 \\
\textbf{nu} & Hz & Frequency of the Chopper, omega=2*PI*nu &  \\
jitter & s & Jitter in the time phase & 0 \\
delay & s & Time 'delay'. & 0 \\
isfirst & 0/1 & Set it to 1 for the first chopper position in a cw source & 0 \\
abs\_out & 0/1 & Absorb neutrons hitting outside of chopper radius? & 1 \\
phase & deg & Angular 'delay' (overrides delay) & 0 \\
verbose & 1 & Set to 1 to display Disk chopper configuration & 0 \\
n\_pulse & 1 & Number of pulses (Only if isfirst) & 1 \\
sequence & str & Slit sequence around disk, with 0=closed, 1=open positions. & "NULL" \\
\bottomrule
\end{longtable}

\subsection*{Links}
\begin{itemize}
  \item Component source code found in file \texttt{StatisticalChopper.comp}.
  \item R. Von Jan and R. Scherm. The statistical chopper for neutron time-of-flight spectroscopy. Nuclear Instruments and Methods, 80 (1970) 69-76.
\end{itemize}
\IfFileExists{contrib/StatisticalChopper_static.tex}{\input{contrib/StatisticalChopper_static.tex}}{}